% NichidaiMasterExam_R6_3rd_4
\begin{tcolorbox} %%R6_3rd_4
	$\fbox{4}$ 次の問いに答えよ。
\begin{enumerate}
	\item 極座標表示$x= r\cos{\theta}, y= r\sin{\theta}$に対して、Jacobian$\frac{\partial(x, y)}{\partial(r, \theta)}$を求めよ。
	\item $\int x^2\cdot 2^{x^3} dx$を求めよ。
	\item $\iint_{D} 2^{(x^2+ y^2)^{3/2}} dxdy$を求めよ。ただし
\begin{equation*}
	D= \{(r\cos{\theta}, r\sin{\theta});\ 0\le \theta\le 1,\ 2\theta\le r\le 2\}
\end{equation*}
\end{enumerate}
\end{tcolorbox}

\begin{souanswer} %%R6_3rd_4_ans
	$\fbox{4}$
\begin{enumerate}
	\item 
\begin{equation*}
	\frac{\partial (x, y)}{\partial (r, \theta)}= \begin{vmatrix}{\partial x}/{\partial r}& {\partial y}/{\partial r}\\ {\partial x}/{\partial \theta}& {\partial y}/{\partial \theta}\end{vmatrix}
\end{equation*}
	それぞれの偏微分の値を求める。
	${\partial x}/{\partial r}= \cos{\theta}$, 
	${\partial y}/{\partial r}= \sin{\theta}$, 
	${\partial x}/{\partial \theta}= -r\sin{\theta}$, 
	${\partial y}/{\partial \theta}= r\cos{\theta}$.\\
	求めるJacobianは$\begin{vmatrix}\cos{\theta}& \sin{\theta}\\ -r\sin{\theta}& r\cos{\theta}\end{vmatrix}= r$

	\item $t= x^3$とおく。$dt= 3x^2 dx$
\begin{align*}
	I
	&= \int 2^t \left(\frac{1}{3}\right) dt\\
	&= \frac{1}{3\log{2}} 2^{x^3}+ C\ \ (C:積分定数)
\end{align*}

	\item 座標変換して$x^2+ y^2= r^2$、(1)よりJacobianは$r$であるから$dxdy= r drd\theta$
\begin{equation*}
	I= \int_{\theta= 0}^1 \int_{r= 2\theta}^2 2^{r^3}\cdot rdrd\theta
\end{equation*}

\begin{tcolorbox}
\begin{itemize}
	\item $r$の範囲: $r$の最小値は$2\theta= 0(\leadsto \theta= 0)$、最大値は2。範囲は$0\le \theta\le 2$
	\item $\theta$の範囲: $2\theta\le r$より$\theta\le \frac{r}{2}$、$0\le \theta$。範囲は$0\le \theta\le \frac{r}{2}$
\end{itemize}
\end{tcolorbox}
\begin{align*}
	I
	&= \int_{r= 0}^2\int_{\theta= 0}^{r/2} r\cdot 2^{r^3}drd\theta\\
	&= \int_{r= 0}^2\left[\theta\cdot r\cdot 2^{r^3}\right]_{\theta= 0}^{r/2}dr\\
	&= \frac{1}{2}\int_{r= 0}^2 r^2\cdot 2^{r^3}dr\\
	&= \frac{1}{2}\left[\frac{1}{3\log{2}}\cdot 2^{r^3}\right]_{r= 0}^2\\
	&= \frac{1}{6\log{2}}\cdot (2^8- 2^0)\\
	&= \frac{85}{2\log{2}}
\end{align*}
\end{enumerate}
\end{souanswer}